\section{Introduction}

The chemical engineering literature has expanded rapidly across diverse subfields, from flow batteries and PFAS degradation to electrocatalysis and reaction engineering \cite{goodenough2013battery_knowledge,wang2021pfas}. Each year, thousands of experimental papers report new material systems, operating conditions, and performance metrics. Yet the corresponding experimental data and mechanistic knowledge remain locked in unstructured text. Cross-system, cross-material integration is increasingly necessary, particularly when researchers seek to transfer insights from one electrochemical system to another or compare degradation pathways across different material classes.

Experimental conclusions in chemical engineering are rarely determined by a single variable. Instead, each finding reflects the co-constraint of multiple entities operating together: the material system under study, the experimental conditions applied, the performance metrics measured, and the mechanistic evidence supporting the interpretation. These entities span multiple length scales. At the molecular and material scale, composition, crystal structure, functional groups, and defect densities govern intrinsic properties. At the transport and reaction scale, momentum, heat, and mass transfer dictate how materials behave under operating conditions. At the process scale, parameters such as temperature, pressure, pH, concentration, flow rate, and current density define the experimental envelope within which observations are made. Together, these variables form what we term \emph{multi-variable co-constrained experimental fact units} --- each reported conclusion in the literature emerges from the interplay of several such variables rather than any one in isolation.

Vector-based retrieval augmented generation (RAG) systems locate semantically similar text chunks through dense embedding similarity \cite{lewis2020rag}. While effective for locating passages that share topical overlap with a query, this approach offers no structural guarantee that all variables belonging to the same experimental fact will be recalled together. A query about a specific material under certain conditions may retrieve a chunk containing the material name and a separate chunk mentioning the conditions, yet the two chunks may originate from different experimental contexts. For complex conditional queries involving multiple co-constrained variables, vector retrieval tends to fragment the very relationships the user seeks to reconstruct.

Graph-RAG methods introduce structured knowledge representations to address some of these limitations \cite{edge2024local_to_global}. By modeling entities as vertices and relationships as edges, binary graphs bring relational structure into the retrieval process. However, a binary edge fundamentally connects only two vertices. A multi-variable experimental fact such as ``Material A tested under Condition B achieves Metric C via Mechanism D'' is necessarily decomposed into pairwise fragments: A--B, B--C, C--D, and so on. In this decomposition, the key semantic information --- that these variables belong to \emph{the same experimental fact} --- is distributed across multiple independent edges and may not be recoverable during retrieval. The joint constraint that defines the experimental conclusion is lost.

Hypergraphs provide a natural modeling framework for relations that extend beyond pairwise coupling. A hypergraph $H = (V, E)$ generalizes the standard graph by allowing each hyperedge $e \in E$ to connect an arbitrary number of vertices, with $|e| \geq 2$ \cite{bretto2013hypergraph}. When $|e| \geq 3$, the hyperedge encodes a higher-order relation among its constituent vertices that cannot be expressed through pairwise edges alone. Feng et al. \cite{feng2025hyperrag} recently proposed Hyper-RAG, a retrieval framework that employs hypergraphs to preserve higher-order relational structures in knowledge-intensive question answering. In the context of chemical engineering literature, hyperedges offer a direct means to represent multi-variable experimental contexts: each hyperedge can encompass the material, conditions, metrics, and mechanistic evidence that jointly constitute a single experimental fact.

We propose Hyper-ChE, a domain-adaptive hypergraph-enhanced retrieval framework designed specifically for chemical engineering experimental literature. The framework organizes experimental knowledge as \emph{Experimental Fact Units} (EFUs), each modeled as a hyperedge connecting all relevant entities within a reported experimental conclusion. We instantiate Hyper-ChE across two representative domains --- redox flow batteries and PFAS removal/degradation --- and conduct subset-scale validation using corpora of approximately 50 documents per domain against a binary Graph-RAG baseline. Our preliminary evidence indicates that Hyper-ChE retrieves richer multi-variable context and preserves the structural integrity of experimental facts more effectively than pairwise decomposition approaches.

The contributions of this work are as follows:
\begin{enumerate}
  \item We identify Experimental Fact Units (EFUs) as a recurring organizational structure in chemical engineering literature, where each experimental conclusion emerges from the co-constraint of multiple variables across material, condition, metric, and mechanism dimensions.
  \item We propose Hyper-ChE, a domain-adaptive hypergraph knowledge modeling and dual-path retrieval framework that represents EFUs as hyperedges and supports both hypergraph and subgraph retrieval strategies tailored to chemical engineering queries.
  \item We instantiate the framework in two distinct application domains --- redox flow batteries and PFAS removal/degradation --- constructing domain-specific entity ontologies and relation schemas that reflect the multi-scale structure of experimental knowledge in each area.
  \item We present subset-scale validation results comparing Hyper-ChE against a binary Graph-RAG baseline, with preliminary evidence suggesting richer multi-variable context retrieval and improved preservation of co-constrained experimental semantics.
  \item We release prototype implementation code and an interactive demonstration to support reproducibility and community extension to additional chemical engineering subfields.
\end{enumerate}

The remainder of this paper is organized as follows. Section~\ref{sec:background} reviews background on hypergraph modeling and domain-adaptive retrieval. Section~\ref{sec:framework} presents the Hyper-ChE framework, including the EFU definition, hypergraph construction pipeline, and dual-path retrieval mechanism. Sections~\ref{sec:case_flow_battery} and~\ref{sec:case_pfas} describe the two case studies with domain-specific instantiations. Section~\ref{sec:discussion} discusses the findings and their implications, and Section~\ref{sec:conclusion} concludes with candidate directions for further validation.
